% TOP_PROBLEM: {"id": 20001056, "title": "Frontier Hall obstructions for exact distance shells", "category_id": 2, "category": {"id": 2, "name": "combinatorics", "display_name": "Combinatorics", "description": "Counting problems, graph theory, discrete structures.", "slug": "combinatorics", "order_index": 2, "created_at": "2026-07-31T15:26:25.671Z"}, "status": "partially_solved", "problem_number": "AIM-COMBINATORICS-0181", "set_id": 14, "statusQualification": "Makes explicit the nontrivial range k>=2, consistent with the source definition of positive-distance neighbourhoods."}
% FIRST_POSTED: 2026-10-01
% ENTRY_KIND: statement_only
% UnsolvedMath ID 20001056; source catalogue code AIM-COMBINATORICS-0181
% !TeX program = lualatex
% Definition and sources only; this file is not an LLM solution attempt.
\documentclass[11pt,a4paper]{article}
\usepackage[margin=25mm]{geometry}
\usepackage{fontspec}
\setmainfont{lmroman10-regular.otf}[BoldFont=lmroman10-bold.otf,
 ItalicFont=lmroman10-italic.otf,BoldItalicFont=lmroman10-bolditalic.otf]
\usepackage{amsmath,amssymb,mathtools}
\usepackage{unicode-math}
\setmathfont{latinmodern-math.otf}
\AtBeginDocument{\let\setminus\smallsetminus}
\usepackage{xurl,hyperref}
\hypersetup{colorlinks=true,urlcolor=blue}
\setcounter{secnumdepth}{0}
\setlength{\parindent}{0pt}
\setlength{\parskip}{5pt}
\setlength{\emergencystretch}{3em}
\begin{document}
\section*{Frontier Hall obstructions for exact distance shells}\label{20001056}
{\small UnsolvedMath ID 20001056 · AIM-COMBINATORICS-0181 · Combinatorics · AIM Workshop Problem Lists}
\subsection{Definitions and mathematical statement}
For a finite simple digraph $D$, define the outgoing distance
$d_D(v,w)$ to be the minimum length of a directed path from $v$ to
$w$, taking it to be infinity when no path exists and setting
$d_D(v,v)=0$. For each positive integer $j$, its $j$th outgoing
layer is
\[
                       N_j^+(v)=\{w:d_D(v,w)=j\}.
\]
This is an exact-distance layer, so its members cannot have a shorter
path from $v$ even if they also occur at the end of a $j$-edge walk.
Different paths ending at the same vertex are not counted separately.

Fix an integer $k\ge2$. Suppose $D$ is nonempty and has no directed
cycle of length at most $k$; in particular, it is loopless and has
no digons. The proposed higher-neighbourhood conjecture asserts that
some vertex $v$ satisfies
\[
                          |N_k^+(v)|\ge|N_{k-1}^+(v)|.
\]
The assertion is required for every $k\ge2$ and every such finite
digraph, with the vertex allowed to depend on both $D$ and $k$.
Layers may be empty, and no minimum-degree assumption is made.
For $k=2$ this is the second-neighbourhood conjecture. Larger $k$
compare two successive layers rather than a layer to the entire
ball preceding it. The range $k\ge2$ also keeps both indices within
the positive-layer notation used by the source.
\subsection{Short English statement}
In a digraph with no cycle of length at most k, must some vertex have a kth outgoing layer at least as large as its preceding layer?
\subsection{Sources}
\begin{itemize}
\item[S1] ULAM.AI and the UnsolvedMath Contributors, supplied problems.json, record 20001056 (AIM-COMBINATORICS-0181), AIM Workshop Problem Lists. Catalogue identity and source transcription; CC BY 4.0.
\url{https://huggingface.co/datasets/ulamai/UnsolvedMath}.
\item[S2] American Institute of Mathematics, The Caccetta-Haggkvist conjecture, item 6.13. Original workshop question underlying AIM-COMBINATORICS-0181.
\url{https://aimath.org/WWN/caccetta/caccetta.pdf}.
\end{itemize}
\end{document}
